\subsection{Definition and universal property of the Clifford algebra.}

DEFINITION 1. --- One calls the Clifford algebra of Q and denotes by C(Q) the quotient algebra of the tensor algebra T(E) of the module E by the two-sided ideal (denoted I(Q)) generated by the elements of the form \(x \otimes x - Q(x)\) . 1 ( \(x \in E\) ).

We will denote by \(\rho_{\mathbb{Q}}\) (or simply \(\rho\) when no confusion is to be feared) the mapping from E into C(Q) composed of the canonical mapping from E into T(E) and the canonical mapping \(\sigma\) of T(E) onto C(Q); the map \(\rho_{\mathbb{Q}}\) is said to be canonical. Note that C(Q) is generated by \(\rho_{\mathbb{Q}}(E)\) , and that, for \(x \in E\) , one has

\[
\rho (x) ^ {2} = \mathrm{Q} (x). 1;\tag{1}
\]

whence, replacing \(x\) by \(x + y\) ( \(x, y\) in E)

\[
\rho (x) \rho (y) + \rho (y) \rho (x) = \Phi (x, y). 1\tag{2}
\]

Example. --- If E admits a basis consisting of a single element \(e\) , T(E) is isomorphic to the polynomial algebra A{[}X{]}, and C(Q) is a quadratic extension of A, with basis \((1, u)\) , where \(u\) is the element \(u = \rho(e)\) and satisfies \(u^2 = Q(e)\) .

Denote by \(T^{h}\) the h-th tensor power \(\otimes E\) in \(\mathrm{T}(\mathrm{E})\) , and let \(T^{+}\) (resp. \(T^{-}\) ) be the sum of \(T^{h}\) for h even (resp. odd).

Since T(E) is the direct sum of \(T^{+}\) and \(T^{-}\) and since I(Q) is generated by elements of \(T^{+}\) , I(Q) is the direct sum of \(T^{+} \cap I(Q)\) and \(T^{-} \cap I(Q)\) , and C(Q) is the direct sum of the two submodules \(C^{+}(Q) = \sigma(T^{+})\) and \(C^{-}(Q) = \sigma(T^{-})\) (also denoted \(C^{+}\) and \(C^{-}\) ). The elements of \(C^{+}\) will be called even (resp. odd). We have the relations

\[
\mathrm{C} ^ {+} \mathrm{C} ^ {+} \subset \mathrm{C} ^ {+}, \quad \mathrm{C} ^ {+} \mathrm{C} ^ {-} \subset \mathrm{C} ^ {-}, \quad \mathrm{C} ^ {-} \mathrm{C} ^ {+} \subset \mathrm{C} ^ {-}, \quad \mathrm{C} ^ {-} \mathrm{C} ^ {-} \subset \mathrm{C} ^ {+}.\tag{3}
\]

In particular \(C^{+}\) is a subalgebra of \(C(Q)\) .

PROPOSITION 1. --- Let f be a linear map of E into an algebra D over A such that \(f(x)^{2} = \mathrm{Q}(x)\) . 1 for every \(x \in \mathbf{E}\) . There exists one and only one homomorphism \(\bar{f}\) of C(Q) into D such that \(f = \bar{f} \circ \rho_{\mathbb{Q}}\) .

The uniqueness of \(\bar{f}\) follows from the fact that C(Q) is generated by \(\rho_{\mathbb{Q}}(\mathrm{E})\) . Let \(h\) be the unique homomorphism of T(E) into D extending \(f\) ( \(h\) is defined by \(h(x_1 \otimes \cdots \otimes x_n) = f(x_1) \ldots f(x_n)\) ). We have

\[
h (x \otimes x - \mathrm{Q} (x). 1) = (f (x) ^ {2} - \mathrm{Q} (x)). 1 = 0,
\]

and consequently h vanishes on I(Q) and defines, by passage to the quotient, the homomorphism \(\bar{f}\) sought.

Prop. 1 expresses that C(Q) is the solution of a universal mapping problem (Ens., chap.~IV, § 3, no. 1).

Take in particular for D the opposite algebra of C(Q) and for \(f\) the map \(\rho\) ; prop. 1 implies that there exists one and only one anti-automorphism \(\beta\) of C(Q) whose restriction to \(\rho(E)\) is the identity; it is called the principal anti-automorphism of C(Q). It is clear that \(\beta^{2} = 1\) .

On the other hand, let Q' be a quadratic form on an A-module E' and f a linear map from E to E' such that Q'∘f = Q. We have \(\rho_{\mathbb{Q}'}(f(x))^{2} = \mathrm{Q}'(f(x)).1 = \mathrm{Q}(x).1\) , and consequently there exists one and only one homomorphism C(f) from C(Q) to C(Q') such that C(f)∘ρq = ρq'∘f.~If f is the identity, C(f) is the identity; if Q'' is a quadratic form on an A-module E'' and g a linear map from E' to E'' such that Q''∘g = Q', we have C(g∘f) = C(g)∘C(f). When E' is a submodule of E and f the canonical injection of E' into E (so that Q' is the restriction of Q to E'), we say that C(f) is the canonical homomorphism of C(Q') into C(Q).

Take in particular Q' = Q and for f the map \(x \to -x\) ; we see that there exists one and only one automorphism \(\alpha\) of C(Q) such that \(\alpha \circ \rho = -\rho\) ; it is called the principal automorphism of C(Q). It is clear that \(\alpha^{2} = 1\) , and that the restriction of \(\alpha\) to C⁺ (resp. C⁻) is the identity (resp. the map \(u \to -u\) ).

PROPOSITION 2. --- Let A' be a commutative ring, φ a homomorphism from A to A', Q' the quadratic form on E' = A' ⊗\_A E deduced from Q by extension of scalars (§ 3, n° 4, prop. 3). There exists one and only one isomorphism j from the algebra A' ⊗\_A C(Q) onto C(Q') such that j (1 ⊗ ρ\_Q(x)) = ρ\_Q′(1 ⊗ x) for all x ∈ E.

It suffices to show that the algebra \(C' = A' \otimes C(Q)\) and the map \(1 \otimes_{\rho_{Q}}\) of \(E'\) into \(C'\) form a solution of the same universal problem as \(C(Q')\) and \(\rho_{Q'}\) . Indeed, let \(D'\) be an algebra over \(A'\) and \(f'\) an A'-linear map from \(E'\) in \(D'\) such that \(f'(x')^{2} = Q'(x')\) . 1 for every \(x' \in E'\) . The map \(g : x \to f'(1 \otimes x)\) from E into \(D'\) (considered as an A-module via the homomorphism \(\varphi\) ) is A-linear and we have \(g(x)^{2} = Q'(1 \otimes x)\) . 1 = Q(x). 1 for all \(x \in E\) . There therefore exists one and only one A-homomorphism \(\bar{g}\) from C(Q) to \(D'\) such that \(\bar{g}(\rho_{Q}(x)) = f'(1 \otimes x)\) . Consequently there exists one and only one A'-homomorphism \(\overline{f'}\) from C' to \(D'\) such that \(\overline{f}'(1 \otimes \rho_{Q}(x)) = f'(1 \otimes x)\) for all \(x \in E\) ; by linearity it follows that \(\overline{f}'((1 \otimes \rho_{Q})(x')) = f'(x')\) for all \(x' \in E'\) . QED.

